\section{Introduction}
Automatic prompt and program optimisation is now a routine step in building LLM systems: a tool such as DSPy/MIPRO \citep{khattab2023dspy,opsahl2024optimizing} or an instruction-search procedure \citep{zhou2022large} evaluates many candidate configurations (instructions, demonstrations, formats, module wirings) on the items of a benchmark and returns the best one. The number that then goes into the paper, the model card or the product decision is, almost always, the winner's score on \emph{those same items}.
That number is a maximum over noisy estimates. Even if every candidate were exactly as good as every other, the best of $B$ of them would look better than average by roughly $\sigma\sqrt{2\ln B}$, where $\sigma$ is the evaluation noise of a single candidate on $n$ items; and decisions made from such numbers --- ``our tuned pipeline beats the baseline'', ``method A beats method B'' --- inherit the bias, particularly when the two sides were tuned with different effort.

This paper asks what a reported score means after automatic tuning, and how to report one that means something.
We make four contributions.
\begin{enumerate}\setlength\itemsep{0pt}
\item \textbf{A quantitative account of the gap.} For exhaustive search over $B$ candidates we give a closed form for the expected optimism, $\sigma^2(\tau^2+\sigma^2)^{-1/2}\,\mathbb{E}Z_{(B)}$, which makes explicit that it grows with the search budget, shrinks with benchmark size, and is \emph{largest when the candidates truly differ least}, i.e.\ when tuning is least useful (Sec.~\ref{sec:method}); we verify it in 375 simulated settings.
\item \textbf{An audit on a real LLM and benchmark.} We cache the per-item correctness of $72$ prompt programs (answer format $\times$ demonstration set $\times$ instruction $\times$ answer cue) for Qwen2.5-1.5B-Instruct on $\RealN$ MMLU-Pro questions and replay four families of search procedures (random search, successive halving, coordinate ascent, a surrogate-guided MIPRO-like search) on random benchmark subsets of 50--800 items, scoring each selected configuration by its exact accuracy over the whole cached population (Sec.~\ref{sec:exp}).
\item \textbf{Consequences for conclusions.} We measure how often a tuned system is declared significantly better than an untuned or lightly tuned competitor when it is not.
\item \textbf{Reporting methods and how to test them.} We compare the status quo with hold-out reporting, cross-fitting, a replay bootstrap that re-runs the \emph{whole search} on resamples, empirical-Bayes shrinkage, conditional-on-selection inference and simultaneous intervals, and Ladder / Thresholdout, using a reliability protocol: known-truth coverage, a no-signal control, robustness to the shape of the candidate distribution, and decision error rates (Sec.~\ref{sec:exp}).
\end{enumerate}
